\title{Group Sequence Policy Optimization}

\begin{abstract}
We present Group Sequence Policy Optimization (GSPO), a reinforcement learning method designed for robust, efficient, and high-quality training of large language models. GSPO diverges from traditional approaches that rely on token-level importance ratios by employing sequence-based likelihood ratios and implementing sequence-level mechanisms for clipping, reward computation, and optimization. Our experiments show that GSPO not only outperforms the GRPO algorithm in both training efficiency and model results, but also enhances stability in Mixture-of-Experts (MoE) RL training scenarios. Additionally, GSPO streamlines RL infrastructure development. These advantages of GSPO have played a central role in the significant performance advances seen in the latest Qwen3 models.
\end{abstract}

\section{Introduction}
\label{sec:intro}

Reinforcement learning (RL) has become a central approach for advancing the scalability of language models \citep{o1,dpsk-r1,qwq32b,qwen3}. Utilizing RL at scale enables these models to address increasingly challenging tasks, including advanced mathematics and programming, by engaging in extended and more complex reasoning.

To extend RL's effectiveness with higher computational resources, it is crucial to ensure that training processes remain stable and resilient. Nevertheless, leading RL algorithms like GRPO \citep{grpo} face significant instability when applied to very large language models, frequently leading to catastrophic and unrecoverable collapse \citep{qwen3, minimax-m1}. Such instability significantly impedes efforts to enhance the capabilities of language models via ongoing RL training.

This work demonstrates that GRPO's instability originates from a fundamental flaw: the erroneous application and breakdown of importance sampling weights in the algorithm. As a result, substantial variance is injected into the training process, increasing in severity with longer responses. This issue is exacerbated by the use of a clipping mechanism, which collectively leads to the deterioration and eventual collapse of the model.

To overcome these central shortcomings, we introduce \textbf{Group Sequence Policy Optimization (GSPO)}, an RL framework tailored for large language model training. GSPO's principal contribution is its theoretically principled approach to defining the importance ratio, derived from sequence likelihood as described by \citet{click}, thus faithfully adhering to the core tenets of importance sampling. Furthermore, GSPO utilizes normalized rewards to represent the advantages among several responses to the same query, thereby facilitating coherent optimization that reflects sequence-level reward assignments.

Through experimental analysis, we find that GSPO markedly outperforms GRPO in terms of training consistency, speed, and effectiveness. Notably, GSPO naturally overcomes the instability issues commonly encountered during the reinforcement learning of expansive Mixture-of-Experts (MoE) models, thus obviating the requirement for intricate stabilization techniques and indicating opportunities to streamline RL systems. These advantages of GSPO have been instrumental in achieving remarkable performance gains in the updated Qwen3 models. We foresee GSPO as a stable and extensible base for ongoing progress in large-scale RL training applied to language models.

\section{Preliminaries}

\paragraph{Notation}
Throughout this study, the autoregressive language model with parameters $\theta$ is regarded as a policy, represented as $\pi_\theta$. We represent a query by $x$ and designate the entire collection of queries as $\mathcal{D}$. For a given query $x$ and its corresponding response $y$, the probability assigned to $y$ by the policy $\pi_\theta$ is expressed as $\pi_\theta (y | x) = \prod_{t=1}^{|y|} \pi_\theta (y_t | x, y_{<t})$, with $|y|$ indicating the total number of tokens in $y$. Each query-response pair $(x, y)$ is evaluated by a verifier $r$, which assigns a reward $r(x, y)$ where $r(x, y) \in [0, 1]$.

\paragraph{Proximal Policy Optimization (PPO)} Using samples generated from the old policy $\pi_{\theta_\text{old}}$, PPO \citep{ppo} constrains the policy update within a proximal region of the old policy through the clipping mechanism. Specifically, PPO employs the following objective for policy optimization (we omit the KL regularization term hereinafter for brevity, as it is not the focus of this paper):

\begin{align}
\mathcal{J}_\text{PPO}(\theta) = \mathbb{E}_{ x \sim \mathcal{D},\, y \sim \pi_{\theta_\text{old}}( \cdot | x) }
\left[ \frac{1}{|y|} \sum_{t=1}^{|y|} 
\min \left( w_{t}(\theta) \widehat{A}_{t},  \, \mathrm{clip} \left( w_{t}(\theta), 1 - {\varepsilon}, 1 + {\varepsilon}\right) \widehat{A}_{t} \right)
\right],
\end{align}

Here, the token importance ratio $y_t$ is expressed as $
w_{t}(\theta) = \frac{ \pi_{\theta} (y_{t} | x, y_{<t}) }{ \pi_{\theta_\text{old}} (y_{t} | x,y_{<t})} $, the advantage estimate $\widehat{A}_{t}$ for $y_t$ is provided by a separate value model, and $\varepsilon$ denotes the clipping threshold for the importance ratios.

A principal obstacle encountered in PPO involves the value model's centrality to the method. In practical settings, the value model is typically as large as the policy model itself, resulting in substantial demands on memory and computation. In addition, the performance of the algorithm is tightly linked to the precision of value predictions. Obtaining an accurate value model is already difficult, but scaling it to accommodate lengthier outputs and more sophisticated tasks makes this problem considerably more acute.

\paragraph{Group Relative Policy Optimization (GRPO)} GRPO \citep{grpo} bypasses the need for the value model by computing the relative advantage of each response within a group of responses to the same query. Specifically, GRPO optimizes the following objective:

\begin{align}
\label{equ:grpo}
\mathcal{J}_\text{GRPO}(\theta) = \mathbb{E}_{ x \sim \mathcal{D},\, \{y_i\}_{i=1}^G \sim \pi_{\theta_\text{old}}( \cdot | x) }
\left[ \frac{1}{G} \sum_{i=1}^{G} \frac{1}{|y_i|} \sum_{t=1}^{|y_i|} 
\min \left( w_{i,t}(\theta) \widehat{A}_{i,t},  \, \mathrm{clip} \left( w_{i,t}(\theta), 1 - {\varepsilon}, 1 + {\varepsilon}\right) \widehat{A}_{i,t} \right)
\right],
\end{align}

where $G$ is the number of generated responses to each query $x$ (i.e., the group size), and the importance ratio $w_{i,t}(\theta)$ and advantage $\widehat{A}_{i,t}$ of token $y_{i,t}$ are:

\begin{align}
    w_{i,t}(\theta)=\frac{ \pi_{\theta} (y_{i,t} | x, y_{i,<t}) }{ \pi_{\theta_\text{old}} (y_{i,t} | x,y_{i,<t})},\quad\ 
    \widehat{A}_{i,t} = \widehat{A}_{i} = \frac{r(x, y_i) - \mathrm{mean} \left( \{ r(x, y_i) \}_{i=1}^G \right) }{ \mathrm{std} \left( \{ r(x, y_i) \}_{i=1}^G \right) },
\end{align}

respectively, where all the tokens in $y_i$ share the same advantage as $\widehat{A}_{i}$.

\section{Motivation}
\label{sec:motivation}

As models increase in size, incorporate greater sparsity—such as with Mixture-of-Experts architectures—and generate longer responses, reinforcement learning benefits from employing large rollout batch sizes to fully leverage available hardware. To enhance sample efficiency, it is common to divide these substantial batches of rollout data into several smaller mini-batches during the process of gradient computation. However, this approach naturally gives rise to an off-policy learning context, since the responses $y$ are drawn from a previous policy $\pi_{\theta_\text{old}}$, not the policy $\pi_\theta$ currently under optimization. This dynamic underpins the importance of the clipping mechanisms found in PPO and GRPO, which are designed to restrict the influence of samples that deviate significantly from the current policy during gradient estimation.

While mechanisms like clipping aim to manage this off-policy discrepancy, we identify a more fundamental issue in GRPO: \textit{its objective is ill-posed}. This problem becomes particularly acute when training large models on long-response tasks, leading to catastrophic model collapse. The ill-posed nature of the GRPO objective stems from a misapplication of importance sampling weights. The principle of importance sampling is to estimate the expectation of a function $f$ under a target distribution $\pi_\text{tar}$ by re-weighting samples drawn from a behavior distribution $\pi_\text{beh}$:

\begin{align}
\mathbb{E}_{z \sim \pi_\text{tar}} \left[ f(z) \right] = \mathbb{E}_{z \sim \pi_\text{beh}} \left[ \frac{ \pi_\text{tar} (z) }{ \pi_\text{beh} (z)} \, f(z) \right].
\end{align}

Fundamentally, effective distributional correction hinges on taking an average across a substantial number of samples ($N \gg 1$) drawn from the behavior policy $\pi_\text{beh}$, so that the importance weight $\frac{ \pi_\text{tar} (z) }{ \pi_\text{beh} (z)}$ compensates for discrepancies between the behavior and target distributions.

Conversely, GRPO incorporates the importance weight $\frac{ \pi_{\theta} (y_{i,t} | x, y_{i,<t}) }{ \pi_{\theta_\text{old}} (y_{i,t} | x,y_{i,<t})}$ at each individual token timestep $t$. As this weight relies on merely a single sampled token $y_{i,t}$ per conditional distribution $\pi_{\theta_\text{old}}( \cdot | x, y_{i, <t})$, it does not adequately adjust for the distribution gap, but rather injects substantial noise into gradient updates. This variance accumulates across lengthy sequences and is further intensified by the gradient clipping procedure. Our empirical analysis reveals that this phenomenon can precipitate an irreversible collapse in model behavior: after such collapse, attempts to restart training—even by reverting to a prior checkpoint and fine-tuning hyperparameters (such as clipping thresholds), increasing the sequence length, or modifying the RL queries—fail to restore performance.

These findings highlight an inherent flaw in the architecture of GRPO. The breakdown of importance weighting at the token level points to a foundational lesson: \textbf{optimization should operate at the same granularity as the reward signal}. Given that rewards are assigned based on the full sequence, applying off-policy corrections on individual tokens is evidently suboptimal. Thus, we are compelled to abandon the token-wise objective and instead pursue optimization and importance weighting at the \textit{sequence level}.

\section{Algorithm}

\subsection{GSPO: Group Sequence Policy Optimization}

Although the token-level importance weight $\frac{ \pi_{\theta} (y_{i,t} | x, y_{i,<t}) }{ \pi_{\theta_\text{old}} (y_{i,t} | x,y_{i,<t})}$ presents challenges in GRPO, our analysis indicates that, specifically for language generation tasks, the \textit{sequence-level} importance weight $\frac{ \pi_{\theta} (y | x) }{ \pi_{\theta_\text{old}} (y | x)}$ possesses clear theoretical significance. This quantity quantifies the discrepancy between a response $y$ generated under $\pi_{\theta_\text{old}} (\cdot | x)$ and that of $\pi_{\theta} (\cdot | x)$, effectively mirroring sequence-level reward considerations. Furthermore, it offers an interpretable metric for evaluating the function of the clipping mechanism.

Based on this straightforward observation, we propose the \textbf{Group Sequence Policy Optimization (GSPO)} algorithm. GSPO employs the following sequence-level optimization objective:

\begin{align}
\mathcal{J}_\text{GSPO} (\theta) =
\mathbb{E}_{ x \sim \mathcal{D},\, \{y_i\}_{i=1}^G \sim \pi_{\theta_\text{old}}( \cdot | x) }
\left[ 
\frac{1}{G} \sum_{i=1}^{G}
\min \left( s_{i}(\theta)  \widehat{A}_{i},  \, \mathrm{clip} \left( s_{i}(\theta), 1 - {\varepsilon}, 1 + {\varepsilon} \right) \widehat{A}_{i} \right) 
\right],
\label{equ:gspo}
\end{align}

where we adopt the group-based advantage estimation:

\begin{align}
\widehat{A}_{i} = \frac{r(x, y_i) - \mathrm{mean} \left( \{ r(x, y_i) \}_{i=1}^G \right) }{ \mathrm{std} \left( \{ r(x, y_i) \}_{i=1}^G \right) },
\end{align}

and define the importance ratio $s_{i}(\theta)$ based on sequence likelihood \citep{click}:

\begin{align}
s_{i}(\theta) = \left( \frac{ \pi_{\theta} (y_i | x) }{ \pi_{\theta_\text{old}} (y_i | x)} \right)^{\frac{1}{|y_i|}}
=
\exp \left( \frac{1}{|y_i|} \sum_{t=1}^{|y_i|} \log \frac{ \pi_{\theta} (y_{i,t} | x, y_{i,<t}) }{ \pi_{\theta_\text{old}} (y_{i,t} | x,y_{i,<t})} \right).
\end{align}

Consequently, GSPO implements clipping at the response level, rather than at the token level, to prevent excessively "off-policy" samples from contributing to the gradient computation. This approach is consistent with both sequence-level reward assignment and optimization objectives. Additionally, length normalization is utilized in $s_{i}(\theta)$ to minimize variance and to keep $s_{i}(\theta)$ within a standardized numeric interval. Without normalization, minor likelihood variations in select tokens could provoke substantial shifts in the sequence-level importance ratio, requiring different clipping thresholds for responses of varying lengths. Furthermore, it should be emphasized that GSPO and earlier methods such as GRPO generally employ clipping bounds of disparate magnitudes, a result of their differing importance ratio formulations.

\subsection{Gradient Analysis}
\label{subsec:gradient}

We can derive the gradient of the GSPO objective as follows (clipping is omitted for brevity):

\begin{align}
\nabla_{\theta} \mathcal{J}_\text{GSPO} (\theta)
=&\ 
\nabla_{\theta} \mathbb{E}_{ x \sim \mathcal{D},\, \{y_i\}_{i=1}^G \sim \pi_{\theta_\text{old}}( \cdot | x) }
\left[ 
\frac{1}{G} \sum_{i=1}^{G} s_{i}(\theta) \widehat{A}_{i}
\right] \\
= &\ 
\mathbb{E}_{ x \sim \mathcal{D},\, \{y_i\}_{i=1}^G \sim \pi_{\theta_\text{old}}( \cdot | x) }
\left[ 
\frac{1}{G} \sum_{i=1}^{G}
s_{i}(\theta) \widehat{A}_{i}
\cdot \nabla_{\theta} \log s_{i}(\theta)
\right] \\
=&\ 
\mathbb{E}_{ x \sim \mathcal{D},\, \{y_i\}_{i=1}^G \sim \pi_{\theta_\text{old}}( \cdot | x) }
\left[ 
\frac{1}{G} \sum_{i=1}^{G} 
\left( \frac{ \pi_{\theta} (y_i | x) }{ \pi_{\theta_\text{old}} (y_i | x)} \right)^{\frac{1}{|y_i|}} \widehat{A}_{i} 
\cdot \frac{1}{|y_i|} \sum_{t=1}^{|y_i|} \nabla_{\theta} \log \pi_{\theta} (y_{i,t} | x, y_{i,<t}) 
\right].
\label{equ:gspo_grad}
\end{align}

For comparison, the gradient of the GRPO objective is as follows (note that $\widehat{A}_{i,t} = \widehat{A}_{i}$):

\begin{align}
\nabla_{\theta} \mathcal{J}_\text{GRPO} (\theta) 
=&\ 
\nabla_{\theta} \mathbb{E}_{ x \sim \mathcal{D},\, \{y_i\}_{i=1}^G \sim \pi_{\theta_\text{old}}( \cdot | x) }
\left[ \frac{1}{G} \sum_{i=1}^{G} \frac{1}{|y_i|} \sum_{t=1}^{|y_i|} 
w_{i,t}(\theta) \widehat{A}_{i,t}
\right] \\
=&\
\mathbb{E}_{ x \sim \mathcal{D},\, \{y_i\}_{i=1}^G \sim \pi_{\theta_\text{old}}( \cdot | x) }
\left[ \frac{1}{G} \sum_{i=1}^{G} \widehat{A}_{i} 
\cdot \frac{1}{|y_i|} \sum_{t=1}^{|y_i|} 
\frac{ \pi_{\theta} (y_{i,t} | x, y_{i,<t}) }{ \pi_{\theta_\text{old}} (y_{i,t} | x,y_{i,<t})} 
\nabla_{\theta} \log \pi_{\theta} (y_{i,t} | x, y_{i,<t})  
\right].
\end{align}

Consequently, the primary difference between GSPO and GRPO is rooted in the \textit{method used to assign weights to the gradients of token log likelihoods}. GRPO employs individualized weights for tokens, determined by the "importance weight" $\frac{ \pi_{\theta} (y_{i,t} | x, y_{i,<t}) }{ \pi_{\theta_\text{old}} (y_{i,t} | x,y_{i,<t})}$ for each one. These weights are uneven and can fluctuate within the ranges $(0, 1+\varepsilon]$ when $\widehat{A}_{i} >0$, or $[1-\varepsilon, +\infty)$ for $\widehat{A}_{i} < 0$. Such disparities are significant, and their effects may compound during training, causing unpredictable behavior. By comparison, GSPO applies uniform weighting across all tokens within a response, thus avoiding the instability inherent to GRPO’s approach.

\subsection{GSPO-token: A Token-level Objective Variant}

In scenarios like multi-turn RL, we may desire a finer-grained advantage adjustment than the sequence level. To this end, we introduce a token-level objective variant of GSPO, namely \textbf{GSPO-token}, to allow token-wise advantage customization:

\begin{align}
\mathcal{J}_\text{GSPO-token}(\theta)
=
\mathbb{E}_{ x \sim \mathcal{D},\, \{y_i\}_{i=1}^G \sim \pi_{\theta_\text{old}}( \cdot | x) }
\left[ 
\frac{1}{G} \sum_{i=1}^{G} \frac{1}{|y_i|} \sum_{t=1}^{|y_i|} 
\min \left( s_{i,t}(\theta) \widehat{A}_{i,t},  \, \mathrm{clip} \left( s_{i,t}(\theta), 1 - {\varepsilon}, 1 + {\varepsilon} \right) \widehat{A}_{i,t} \right) 
\right],
\label{equ:gspo-token}
\end{align}

where

\begin{align}
s_{i,t}(\theta) 
= \mathrm{sg} \left[ s_{i}(\theta) \right]  \cdot \frac{ \pi_{\theta} (y_{i,t} | x, y_{i,<t}) }{ \mathrm{sg} \left[ \pi_{\theta} (y_{i,t} | x, y_{i,<t}) \right] },
\end{align}

and $\mathrm{sg}[\cdot]$ denotes only taking the numerical value but stopping the gradient, corresponding to the \texttt{detach} operation in PyTorch. The gradient of GSPO-token can be derived as:

\begin{align}
\nabla_{\theta} \mathcal{J}_\text{GSPO-token}(\theta)
=&\ 
\nabla_{\theta} \mathbb{E}_{ x \sim \mathcal{D},\, \{y_i\}_{i=1}^G \sim \pi_{\theta_\text{old}}( \cdot | x) }
\left[ 
\frac{1}{G} \sum_{i=1}^{G}
\frac{1}{|y_i|} \sum_{t=1}^{|y_i|} 
s_{i,t}(\theta) \widehat{A}_{i,t}
\right] \\
=&\ 
\mathbb{E}_{ x \sim \mathcal{D},\, \{y_i\}_{i=1}^G \sim \pi_{\theta_\text{old}}( \cdot | x) }
\left[ 
\frac{1}{G} \sum_{i=1}^{G} s_i (\theta) \cdot \frac{1}{|y_i|} \sum_{t=1}^{|y_i|} 
\widehat{A}_{i,t} \frac{ \nabla_{\theta} \pi_{\theta} (y_{i,t} | x, y_{i,<t}) }{ \pi_{\theta} (y_{i,t} | x, y_{i,<t}) }
\right] \\
=&\ 
\mathbb{E}_{ x \sim \mathcal{D},\, \{y_i\}_{i=1}^G \sim \pi_{\theta_\text{old}}( \cdot | x) }
\left[ 
\frac{1}{G} \sum_{i=1}^{G}  \left( \frac{ \pi_{\theta} (y_i | x) }{ \pi_{\theta_\text{old}} (y_i | x)} \right)^{\frac{1}{|y_i|}}
\cdot \frac{1}{|y_i|} \sum_{t=1}^{|y_i|}
\widehat{A}_{i,t} \nabla_{\theta} \log \pi_{\theta} (y_{i,t} | x, y_{i,<t})  
\right].
\label{equ:gspo-token_grad}
\end{align}

Observe that $\frac{ \pi_{\theta} (y_{i,t} | x, y_{i,<t}) }{ \mathrm{sg} \left[ \pi_{\theta} (y_{i,t} | x, y_{i,<t}) \right] }$ evaluates numerically to 1. Consequently, $s_{i,t}(\theta)$ is equivalent in value to $s_{i}(\theta)$. Examining Equation~\eqref{equ:gspo} alongside Equation~\eqref{equ:gspo-token}, as well as Equation~\eqref{equ:gspo_grad} and Equation~\eqref{equ:gspo-token_grad}, it is apparent that GSPO-token and GSPO yield the same numerical results regarding the optimization objective, the clipping criterion, and the theoretical gradient, provided that each token in the response $y_i$ is assigned a uniform advantage, specifically $\widehat{A}_{i,t} = \widehat{A}_{i}$. Nonetheless, GSPO-token confers greater adaptability, enabling the assignment of distinct advantages to individual tokens.

\section{Experiments and Discussion}

\subsection{Empirical Results}

Our experiments utilize a cold-start model that has been fine-tuned from Qwen3-30B-A3B-Base. We present both the training reward trajectories and model performance over time on several benchmarks: AIME'24 (reporting average Pass@1 across 32 samples), LiveCodeBench (202410-202502, averaging Pass@1 over 8 samples), and CodeForces (measured via Elo Rating). For reinforcement learning, the rollout data collected in each batch is divided into four mini-batches, which are then used for computing gradient updates. Within the GSPO framework, we specify the left and right clipping bounds in Equation~\eqref{equ:gspo} as 3e-4 and 4e-4, respectively. Our baseline comparison is conducted using GRPO, where we assign the left and right clipping thresholds in Equation~\eqref{equ:grpo} to 0.2 and 0.27, ensuring parity by thorough hyperparameter tuning. It is important to highlight that GRPO relies on the Routing Replay approach for proper convergence of MoE RL (further explored in \S~\ref{subsec:routing_replay}), whereas \textit{GSPO eliminates the necessity for this strategy}.

As depicted in Figure~\ref{fig:results}, GSPO facilitates steady training progress. Our findings indicate that \textbf{GSPO enables sustained advancements in model performance by scaling up computational resources, frequently refreshing the query set, and allowing for longer generation sequences}. Additionally, GSPO outperforms GRPO with respect to training efficiency, attaining higher accuracy and benchmark results given identical computational budgets and numbers of queries. Furthermore, deploying GSPO within RL training for the newest Qwen3 models underscores its effectiveness, highlighting GSPO’s capacity to harness RL scaling benefits in large language model development.

\subsection{Curious Observation on Clipping Fractions}

An important differentiator between GSPO and GRPO lies in their respective clipping strategies: GSPO applies clipping to whole responses, whereas GRPO targets individual tokens. As illustrated in Figure~\ref{fig:clipping}, this results in GSPO and GRPO exhibiting a difference of two orders of magnitude in the proportion of clipped tokens, and even when the clipping ranges are modified, the magnitude discrepancy remains constant. Notably, although GSPO clips a substantially greater number of tokens, resulting in fewer tokens available for training or gradient estimation, it nonetheless demonstrates higher training efficiency compared to GRPO. This seemingly paradoxical result—where a larger fraction of tokens being clipped corresponds to improved training efficiency—suggests that GRPO's token-level gradient estimates suffer from intrinsic noise and are less effective for exploiting samples. In contrast, the sequence-level methodology employed by GSPO yields a more robust and informative learning signal.

\subsection{Benefit of GSPO for MoE Training}
\label{subsec:routing_replay}

\paragraph{Background}
Unlike dense models, which typically activate all parameters during training, MoE models employ sparse activation, leading to distinct stability concerns during RL training. Notably, when applying the GRPO algorithm, we observed that the \textit{expert-activation volatility} inherent to MoE models poses obstacles to effective RL convergence. Specifically, after one or several gradient update steps, there can be a substantial shift in the set of experts engaged for identical responses. As an illustration, with the 48-layer Qwen3-30B-A3B-Base architecture, it is common to see around 10\% discrepancy between the experts activated under a freshly updated policy $\pi_{\theta}$ versus the previous policy $\pi_{\theta_\text{old}}$ for the same rollout, following each RL gradient update. 

This dynamic, which intensifies in deeper MoE configurations, causes significant oscillations in the token-level importance weights $w_{i,t}(\theta) = \frac{ \pi_{\theta} (y_{i,t} | x, y_{i,<t}) }{ \pi_{\theta_\text{old}} (y_{i,t} | x,y_{i,<t})}$. As outlined in \S~\ref{sec:motivation} and~\ref{subsec:gradient}, this volatility undermines the validity of these importance ratios, thereby obstructing reliable RL convergence.

\paragraph{Our Previous Approach} In earlier work, we addressed this issue by leveraging the \textbf{Routing Replay} training method. This involved recording the experts that were activated by $\pi_{\theta_\text{old}}$, and then reproducing these specific routing selections under $\pi_{\theta}$ during computation of the importance ratios: $w_{i,t}(\theta) = \frac{ \pi_{\theta} (y_{i,t} | x, y_{i,<t}) }{ \pi_{\theta_\text{old}} (y_{i,t} | x,y_{i,<t})}$. Consequently, both $\pi_{\theta} (y_{i,t} | x, y_{i,<t})$ and $\pi_{\theta_\text{old}} (y_{i,t} | x,y_{i,<t})$ operate through identical expert pathways for each token $y_{i,t}$, thereby maintaining consistency in the token-level importance ratios and guaranteeing that gradient-based updates optimize a stable activated subnetwork. The critical role of Routing Replay in facilitating proper convergence during GRPO training of MoE architectures is depicted in Figure~\ref{fig:routing_replay}.

\paragraph{Benefit of GSPO} While Routing Replay assists in the convergence of GRPO training for MoE models, it does so at the expense of increased memory and communication demands, and it may hinder the effective utilization of the MoE model’s full capacity. GSPO, by contrast, as depicted in Figure~\ref{fig:results}, avoids reliance on Routing Replay altogether. GSPO computes the importance ratios $s_i(\theta)$ through standard methods, allowing for typical convergence behavior and stable optimization dynamics. The crucial realization is that GSPO is concerned solely with sequence likelihood (specifically, $\pi_{\theta} (y_i | x)$), rather than the probability of individual tokens (that is, $\pi_{\theta} (y_{i,t} | x, y_{i,<t})$). Because the MoE model consistently retains its abilities in language modeling, fluctuations in sequence likelihood remain minimal. Therefore, GSPO addresses the problem of instability in expert activation intrinsic to MoE models directly, thereby eliminating the necessity for elaborate strategies such as Routing Replay. The result is a more straightforward and robust training regime, which permits unrestricted use of the model's full representational power.

\subsection{Benefit of GSPO for RL Infrastructure}

Due to the differences in numerical precision between training engines such as Megatron and inference engines like SGLang and vLLM, it is standard practice to recalculate the likelihoods of sampled outputs under the prior policy $\pi_{\theta_\text{old}}$ using the training engine. In contrast, GSPO relies on sequence-level likelihoods rather than those computed at each token, making it considerably more robust to precision inconsistencies. As a result, GSPO allows for the direct utilization of likelihoods produced by the inference engine during optimization, thus eliminating the necessity of recomputation with the training engine. This feature proves particularly advantageous in settings involving partial rollout, multi-turn RL, and architectures where training and inference are separated.

\section{Conclusion}

Group Sequence Policy Optimization (GSPO) is introduced as an innovative reinforcement learning algorithm specifically designed for training large language models. GSPO leverages importance sampling by constructing importance ratios from sequence probabilities and applies sequence-level mechanisms for clipping, reward assignment, and optimization. Empirical results show GSPO offers substantially improved training stability, efficiency, and overall performance relative to GRPO, proving particularly effective for large-scale reinforcement learning with MoE models. This has enabled notable advancements observed in recent iterations of Qwen3 models. With GSPO serving as a robust, scalable methodological foundation, continued expansion of reinforcement learning is anticipated to drive significant progress in artificial intelligence.

\end{document}